\chapter{Technical Implementation}
\label{ch:implementation}

% Chapter introduction
This chapter details the software engineering aspects and practical source code implementation of the SENTINEL architecture. Shifting from abstract theory, the content focuses on the project directory organization, core class structures (Classes/Methods), Docker service configurations, and the design of the Graphical User Interface (Streamlit Dashboard) which facilitates Human-in-the-Loop (HITL) interaction.

\section{Technology Stack and Source Code Structure}
\begin{itemize}
    \item \textbf{Core Technologies:} Python 3.10+, LangGraph, LangChain, FAISS, SQLite, Streamlit, Docker, and llama.cpp (via the Llama-cpp-python binding).
    \item \textbf{Directory Structure:} A detailed tree representation strictly separating components: \texttt{src/core/} (Tier 1 \& 2), \texttt{src/rag/} (Vector DB), \texttt{src/security/} (Guardrails), \texttt{src/dashboard/} (UI), \texttt{experiments/}, and \texttt{tests/}.
\end{itemize}

\section{Tier-1 Implementation: Filtration Engine \& Welford}
\begin{itemize}
    \item \textbf{\texttt{RunningStats} and \texttt{SessionBaseline} Classes:} Python source code implementation for Welford's algorithm. Resolving dictionary memory allocation challenges via IP hashing and cache garbage collection (TTL cache).
    \item \textbf{\texttt{RuleEngine} Class:} Pattern matching design based on Regular Expressions and Static Thresholds.
\end{itemize}

\section{Tier-2 Implementation: AI Agent and RAG}
\begin{itemize}
    \item \textbf{\texttt{OpenAILLMClient} / Llama.cpp Wrapper:} Parameter configurations for the Local LLM (enforcing \texttt{temperature=0.0} to optimize determinism, context size bounding, and \texttt{n\_gpu\_layers} mapping).
    \item \textbf{LangGraph Workflow Setup:} Utilizing \texttt{StateGraph} to define \texttt{analyze\_node}, \texttt{retrieve\_node}, and the conditional routing function \texttt{should\_retrieve}.
    \item \textbf{\texttt{DualRetriever} Class:} Implementing the FAISS Index and rank\_bm25 to achieve Reciprocal Rank Fusion (RRF). Source code overview of the scoring and sorting algorithms.
\end{itemize}

\section{Guardrail and HMAC System Implementation}
\begin{itemize}
    \item \textbf{\texttt{DelimitedDataEncapsulator} Class:} Logic for generating secure hexadecimal tokens via \texttt{os.urandom} and \texttt{hashlib} to wrap log payloads.
    \item \textbf{\texttt{HMACHashing} Class (Database):} Specification of the SQLite database schema (\texttt{threat\_memory.db}) and the trigger functions calculating SHA-256 hashes derived from Previous Hash + Current Data.
\end{itemize}

\section{SOC Administration Dashboard Implementation (Streamlit UI)}
\begin{itemize}
    \item \textbf{Dashboard Architecture:} Application of a Frontend (Streamlit) communicating dynamically with the Backend (SQLite/File logs) through reactive UI components.
    \item \textbf{Display Modules:}
        \begin{itemize}
            \item \textit{Real-time Queue Dashboard:} Visualizing Tier-1 alert metrics and Tier-2 verdicts.
            \item \textit{Threat Memory Analytics:} Displaying APT alert graphs across IP timelines.
            \item \textit{Audit Trail Integrity:} An automated module verifying HMAC Log Chaining, utilizing a traffic-light interface (Green = Valid, Red = Chain Broken).
        \end{itemize}
    \item \textbf{Human-in-the-loop (HITL):} The interface for reviewing (Approve/Reject) AI-recommended verdicts to update actual Firewall Rules.
\end{itemize}

\section{Containerization and Infrastructure Deployment}
\begin{itemize}
    \item \textbf{\texttt{Dockerfile} Specification:} Multi-stage build mechanisms to minimize image size and compile C++ libraries (llama.cpp) with CUDA acceleration support.
    \item \textbf{\texttt{docker-compose.yml} Service Management:} Linking containers for the Dashboard and Agent Engine, alongside shared log volume mounts.
\end{itemize}

In summary, Chapter 4 systematizes the software development process of SENTINEL, bridging core algorithms with user interfaces and virtualized infrastructure. Chapter 5 will subsequently establish the testing environment to quantitatively evaluate these components.
